\section{Related Work}

\subsection{Vision-Language Embodied Planning and Evaluation}

Large language and vision-language models have increasingly been used as high-level planners for embodied agents. SayCan grounds language-model proposals with learned affordance values, thereby favoring actions that are both semantically relevant and executable \cite{ichterCanNotSay2023}. Inner Monologue instead closes the loop through language-form feedback derived from success detectors, scene descriptions, and human corrections \cite{huangInnerMonologueEmbodied2023}. Grounded Decoding incorporates grounded models directly into language generation to improve the realizability of generated plans \cite{huangGroundedDecodingGuiding2023}, while EmbodiedGPT studies embodied planning and control through multimodal chain-of-thought and instruction tuning \cite{muEmbodiedGPTVisionLanguagePreTraining2023}. Together, these methods show that semantic planning must be constrained by the current environment rather than performed from language priors alone.

Evaluation has likewise progressed from text-centric household simulators such as ALFWorld \cite{shridharALFWorldAligningText2021} and visually grounded instruction-following environments such as ALFRED \cite{shridharALFREDBenchmarkInterpreting2020} and TEACh \cite{padmakumarTEAChTaskDrivenEmbodied2022} to broader multimodal suites. EmbodiedBench evaluates vision-driven agents across embodied navigation and manipulation environments, with subsets targeting visual grounding, spatial reasoning, instruction following, and long-horizon execution \cite{yangEmbodiedBenchComprehensiveBenchmarking2025}. We adopt EmbodiedBench because its heterogeneous task categories expose failures that may share the same terminal reward but arise from different intermediate causes. In contrast to work that primarily improves the planner through parameter updates or more capable decoding, VISTA-Skill keeps the executor frozen and studies how interaction evidence should update its external procedural knowledge.

\subsection{Experience Memory and Procedural Skill Evolution}

External memory provides a training-free route for agents to improve from interaction. Reflexion stores verbal feedback in episodic memory and conditions later trials on self-reflection \cite{shinnReflexionLanguageAgents2023a}, while ExpeL extracts reusable natural-language insights and retrieves relevant prior experience without fine-tuning \cite{zhaoExpeLLLMAgents2024}. Voyager grows an executable code-skill library from environment feedback and self-verification during open-ended exploration \cite{wangVoyagerOpenEndedEmbodied2024}. More structured approaches abstract recurring behavior rather than retaining complete trajectories: Agent Workflow Memory induces reusable workflows in offline or streaming settings \cite{wangAgentWorkflowMemory2025}, and ICAL converts imperfect multimodal demonstrations into retrieval-ready programs annotated with causal relations, state changes, subgoals, and task-relevant visual elements \cite{sarchVLMAgentsGenerate2024a}. These methods establish that compact procedural abstractions can be more reusable than raw episode histories, but they generally infer what to store from trajectory-level success, critique, or supervision.

Recent work treats the skill artifact itself as adaptive state. Skill-Pro represents a skill through activation, execution, and termination conditions, aggregates semantic gradients from experience, and uses a non-parametric PPO-style gate to select candidate revisions \cite{miSkillProLearningReusable2026}. EmbodiSkill, a concurrent embodied-agent preprint, distinguishes skill-changing evidence from executor lapses and performs targeted revision rather than indiscriminate rewriting \cite{juEmbodiSkillSkillAwareReflection2026}. Concurrent SkillOpt work casts a skill document as a text-space optimization variable, applies bounded edits, and accepts them only after held-out validation \cite{yangSkillOptExecutiveStrategy2026}. XSkill jointly accumulates action-level experiences and task-level skills from visually grounded multimodal trajectories \cite{jiangXSkillContinualLearning2026}. These works are closest to ours because they retain a frozen executor and evolve reusable knowledge. Their principal decision, however, is made from a completed trajectory or a scored rollout. VISTA-Skill addresses the preceding question at action granularity: whether a mismatch should revise the transient visual belief, persistent action-effect knowledge, a specific field of a reusable skill, or no persistent memory.

\subsection{Structured Visual State for Embodied Agents}

Scene graphs and object-centric maps provide explicit state for grounding long-horizon decisions. ConceptGraphs constructs open-vocabulary 3D scene graphs that support downstream language queries and planning \cite{guConceptGraphsOpenVocabulary3D2024}, while HOV-SG organizes open-vocabulary scene information hierarchically for language-grounded navigation \cite{werbyHierarchicalOpenVocabulary3D2024}. SayPlan retrieves task-relevant subgraphs from large 3D environments and uses simulator feedback to refine infeasible plans \cite{ranaSayPlanGroundingLarge2023}. SG-Nav builds an online 3D scene graph for hierarchical reasoning and invokes re-perception to correct navigation-relevant perception errors \cite{yinSGNavOnline3D2024}. More recent methods retain richer interaction state: MomaGraph represents objects, functional parts, relations, and state changes in a unified scene graph for graph-then-plan embodied reasoning \cite{juMomaGraphStateAwareUnified2026}, and MSGNav preserves multimodal visual evidence in a 3D scene graph for closed-loop zero-shot navigation \cite{huangMSGNavUnleashingPower2026}.

These representations mainly serve perception, navigation, planning, or feasibility checking. Even when the graph is updated after interaction, it is usually consumed as planner context rather than used to regulate the lifecycle of persistent procedural skills. VISTA-Skill instead treats the graph as an episode-level belief under partial observability and serializes only a task-relevant local graph for execution. More importantly, the same predicate vocabulary connects visual evidence to skill activation, action preconditions, expected effects, and termination. This shared interface makes a procedure's predicted transition comparable with an independently constructed evidence transition without treating the current view as the complete world state.

\subsection{Execution Verification and Reliable Memory Update}

Closed-loop embodied systems verify actions or plans at several stages. Affordance grounding down-ranks infeasible proposals before execution \cite{ichterCanNotSay2023}; environment and simulator feedback enable replanning after a deviation \cite{huangInnerMonologueEmbodied2023,ranaSayPlanGroundingLarge2023}; and re-perception can repair an uncertain scene representation \cite{yinSGNavOnline3D2024}. Skill-evolution methods provide a complementary safeguard at the memory level: Skill-Pro scores candidate skills before pool insertion \cite{miSkillProLearningReusable2026}, EmbodiSkill preserves valid guidance when a failure is attributed to an execution lapse \cite{juEmbodiSkillSkillAwareReflection2026}, and SkillOpt rejects edits that do not improve held-out performance \cite{yangSkillOptExecutiveStrategy2026}. Candidate validation is necessary, as a plausible reflection can still introduce negative transfer, but it does not identify which internal memory produced the original mismatch.

To our knowledge, existing lines of work do not jointly provide (i) an action-level discrepancy between a skill-predicted effect and independently extracted visual evidence, (ii) an explicit update target spanning episodic belief, persistent action knowledge, reusable skill, and no update, and (iii) field-level skill revision conditioned on that attribution. VISTA-Skill combines these elements as \emph{visual transition credit assignment}. This formulation is complementary to downstream skill optimizers and validation gates: it controls which evidence is allowed to reach them and which persistent component, if any, should absorb the update.
